\documentclass[11pt]{article}

\usepackage[margin=1in]{geometry}
\usepackage[T1]{fontenc}
\usepackage[utf8]{inputenc}
\usepackage{lmodern}
\usepackage{amsmath,amssymb}
\usepackage{booktabs}
\usepackage{array}
\usepackage{enumitem}
\usepackage{longtable}
\usepackage{xcolor}
\usepackage{hyperref}
\usepackage{listings}

\hypersetup{
    colorlinks=true,
    linkcolor=blue,
    urlcolor=blue
}

\lstset{
    basicstyle=\ttfamily\small,
    breaklines=true,
    columns=fullflexible,
    keepspaces=true,
    frame=single
}

\title{Solver Behavior and Result Semantics}
\author{}
\date{}

\begin{document}
\maketitle

\section{Global Solver Toggle}

The project has one global solver switch in \texttt{master\_config.py}:
\begin{lstlisting}
enhanced_CBS = True  # or False
\end{lstlisting}

That toggle applies to the \emph{entire} project rather than to individual branches:
\begin{itemize}[leftmargin=2em]
    \item \texttt{False} selects vanilla CBS,
    \item \texttt{True} selects Enhanced CBS (ECBS).
\end{itemize}

So the experiment still compares \textbf{classical mapping vs.\ cyclic mapping} on matched run configurations, but both sides use the same selected high-level solver family.

\section{Static Solver Path}

Static runs go through:
\begin{itemize}[leftmargin=2em]
    \item \texttt{dev/mapf/cbs\_solver.py},
    \item \texttt{dev/mapf/mapf\_low\_level\_astar.py},
    \item \texttt{dev/mapf/low\_level\_guidance.py} when the newer low-level guidance toggles are enabled.
\end{itemize}

When \texttt{enhanced\_CBS=False}, the static branch uses standard CBS with admissible low-level A*.
When \texttt{enhanced\_CBS=True}, the same entry point switches to an ECBS-style search that uses:
\begin{itemize}[leftmargin=2em]
    \item a high-level focal-style selection rule,
    \item conflict count as the secondary node preference inside the focal set,
    \item weighted low-level A* replanning.
\end{itemize}

The ECBS implementation no longer uses one hardcoded project-wide bound. Instead, when \texttt{enhanced\_CBS=True}, the active branch contributes its own \texttt{ECBS\_suboptimality} value from \texttt{master\_config.py}. That same value is used for the static or dynamic ECBS low-level weighted A* calls, the focal-bound calculation, and the recorded solver metadata.

\section{Dynamic Solver Path}

Dynamic runs go through:
\begin{itemize}[leftmargin=2em]
    \item \texttt{dev/mapf/time\_expanded\_cbs.py},
    \item \texttt{dev/mapf/time\_expanded\_astar.py},
    \item \texttt{dev/mapf/low\_level\_guidance.py} when the newer low-level guidance toggles are enabled.
\end{itemize}

The search state is:
\[
(\text{position},\text{time}).
\]
At time $t$, the active movement graph is read from
\[
\texttt{mapped\_loop}[t \bmod L],
\]
where $L$ is the loop length.

The same global toggle is respected here as well, so dynamic branches also switch together between vanilla CBS and ECBS rather than maintaining a separate branch-specific solver choice. The selected branch's \texttt{ECBS\_suboptimality} value is passed into the time-expanded ECBS path in the same way as the static solver.

\section{Branch-Level Low-Level Guidance Toggles}

Each branch now also stores two low-level-planning booleans in \texttt{master\_config.py}:
\begin{lstlisting}
true_static_shortest_path_distance = ...
tight_time_horizon = ...
\end{lstlisting}

These are branch-local rather than global.

\subsection*{\texttt{true\_static\_shortest\_path\_distance}}
When this is \texttt{False}, the project keeps the older Manhattan-style low-level heuristic.
When this is \texttt{True}, the low-level planner uses the exact shortest-path distance on the branch's static movement graph as the heuristic.

For static branches, that graph is the composite map itself.
For dynamic branches, the helper code derives a static graph from the shared mapped loop while ignoring time-varying obstacle occupancy. In practice, this means the heuristic follows the static corridor structure rather than merely geometric proximity.

\subsection*{\texttt{tight\_time\_horizon}}
When this is \texttt{False}, the project keeps the older, more conservative finite horizon rule.
When this is \texttt{True}, the low-level planner uses a tighter distance-aware cap based on the static start-to-goal distance, the latest relevant constraint time, and a bounded slack allowance.

This change does not make the search infinite or unconstrained. It simply reduces the amount of gratuitous waiting-time exploration that the older horizon could allow.

\section{Conflict Model}

Both the static and dynamic solver families detect:
\begin{itemize}[leftmargin=2em]
    \item \textbf{vertex conflicts}: two agents occupying the same vertex at the same time,
    \item \textbf{edge conflicts}: two agents swapping across the same connection at the same time.
\end{itemize}

\section{Disappearing-Agent Model}

Once an agent reaches its final path position, it is no longer considered present for later conflict checks. This makes the system behave like a disappearing-agent MAPF model rather than a stay-at-goal blocking model.

This still matters in clustered-target branches because agents reaching the target cluster disappear instead of continuing to block that region for everyone behind them, regardless of whether the active cluster style is compact or spaced.

\section{Why the New Low-Level Controls Exist}

The practical difficulty in the harder branches is not only high-level conflict branching. Some runs spend substantial time in low-level replanning or even in root-path construction. The new branch-level guidance toggles therefore target two specific costs:
\begin{itemize}[leftmargin=2em]
    \item a weak low-level heuristic in corridor-heavy maps,
    \item an overly generous time horizon in dynamic time-expanded searches.
\end{itemize}

They do not change the experiment from CBS/ECBS into another solver family. They change how the existing low-level search is guided.

\section{Runtime and Progress Reporting}

The solver callbacks report elapsed time in roughly five-second increments. This is mainly for visibility during long runs. The relevant constant is currently defined in \texttt{experiments/study/constants.py}.

\section{Solver Status vs.\ Experimental Category}

The low-level solver status is \emph{not} the same thing as the experimental reporting category. The current mapping is:
\begin{center}
\begin{tabular}{ll}
\toprule
solver status & experimental category \\
\midrule
\texttt{solved} & \texttt{successful} \\
\texttt{bad\_setup\_timeout} & \texttt{unfinished} \\
\texttt{no\_solution} & \texttt{unsolvable} \\
exception/setup failure & \texttt{setup\_failed} \\
\bottomrule
\end{tabular}
\end{center}

\section{Why \texttt{bad\_setup\_timeout} Becomes \texttt{unfinished}}

The project interprets a runtime cap as ``solver halted before completing useful search'' rather than as a proof that the instance is unsolvable. Therefore, timeout-style failures are retained as \texttt{unfinished} and can still count toward the paired quota.

\section{Recorded Per-Run Metrics}

When available, a mapping record stores:
\begin{itemize}[leftmargin=2em]
    \item solver family metadata (CBS vs.\ ECBS and the current suboptimality factor when relevant),
    \item halted computation time,
    \item number of detected conflicts at halt,
    \item average path length,
    \item number of high-level nodes expanded,
    \item run identifiers and seed metadata.
\end{itemize}

Average path length is computed only when solution paths exist. It is derived from total solution cost divided by the number of agents.

\end{document}
